\chapter*{Abstract}
\phantomsection
\addcontentsline{toc}{chapter}{Abstract}

\textbf{This final-year project report presents the design and implementation of an electronic services platform} intended to streamline interactions between users and online administrative or organizational procedures. The work was conducted as part of an end-of-study project at EMSI Rabat and included a professional internship at \textbf{Agence de Développement du Digital (ADD)}. \textbf{It follows a structured software engineering approach}, including company context, problem framing, requirements specification, analysis of existing solutions, and a comprehensive needs analysis covering both functional and non-functional aspects.

\textbf{The proposed solution adopts a layered technical architecture} with a RESTful Laravel backend, a React (Vite) frontend, a PostgreSQL database, Laravel Sanctum SPA authentication (session cookies and CSRF protection), and containerized or classical deployment practices. UML viewpoints are used descriptively to clarify responsibilities, interactions, and deployment topology without relying on graphical screenshots in the written deliverable.

\textbf{Implementation details emphasize modularity, validation, error handling, observability hooks, and maintainability.} Testing activities combine unit, integration, and targeted performance checks. \textbf{Deployment considerations address environment separation, configuration management, and operational readiness.}

\textbf{The outcomes confirm technical feasibility} and highlight practical lessons on security trade-offs, data modeling, and continuous delivery. Future work may extend interoperability, workflow automation, analytics, and hardened production governance.

\vspace{0.8cm}
\noindent\textbf{Keywords:} e-services, REST architecture, Laravel, React, PostgreSQL, SPA authentication, security, containerization, testing, deployment.

\cleardoublepage